还需要完善的功能：
√1. 我刚给projects表添加了一个项目类型project_type字段，需要在add_project.html页面中添加一个下拉框选择项目类型（prj_type）字段。
√2. 用户登录时也运行调用存储函数，并将projects表中的所有project_id传递给存储函数全部运行一遍。
√3. 将价格数据存储到数据库中，而且价格数据要与projects表创建外键关系，以便查询时直接通过projects表查询价格数据,并提供查询价格数据的前端页面。
4. 开发合约规划模块，提供用户通过“目标版本：A-B-C-D，C-E-F-G”的方式录入合同规划，系统可通过该路径自动从对应版本的目标成本计算规划签约金额。
5. 开发合同管理模块，包括合同信息的录入、查询、删除、修改等功能。并且合同要与价格数据关联，以便查询时直接通过合同查询价格数据。
6. 开发技术指标数据录入与查询模块。
7. 开发目标成本测算时，根据项目特征、楼宇特征，匹配参考历史项目和楼宇。用户确认要参考历史项目和楼宇后，系统可自动计算
   参考历史项目和楼宇历史造价指标的平均值，计算目标成本。
8.参考造价文件分部分项和措施项目划分标准完善项目信息和楼宇信息表及前后端代码修改
9.历史项目主要材料价格录入与查询模块
√10.完善：query_data页查询项目信息页面显示出mysql数据库的projects表中项目用地面积（project_land_area）、计容面积（project_count_area)、项目类型（prj_type)
11.根据目标成本测算新的表头修改sql语句中的触发器和存储过程函数
12.完善合同交底模块中个别录入框高度太低的问题
√13.完善目标成本查询模块，自动给每个成本科目加上科目编码，编码的格式要求：一级科目以为“1.”、“2.”、“3.”表示，“1.”表示第一个一级科目，“2.”表示第二个一级科目，以此类推。
二级科目用“1.1.”、“1.2.”、“1.3.”表示，“1.1.”表示第一个二级科目，“1.2.”表示第二个二级科目，以此类推。
三级科目用“1.1.1.”、“1.1.2.”、“1.1.3.”表示，“1.1.1.”表示第一个三级科目，“1.1.2.”表示第二个三级科目，以此类推。如果有必要，继续编写第四级科目、第五级科目等。